\documentclass[aps,prx,reprint,nofootinbib,superscriptaddress]{revtex4-2}

\usepackage{graphicx}
\usepackage{hyperref}
\usepackage{amsmath, amssymb, amsthm}
\usepackage{booktabs}
\usepackage{siunitx}

\newtheorem{theorem}{Theorem}
\newtheorem{lemma}[theorem]{Lemma}
\newtheorem{definition}{Definition}

% Auto-generated by update_counts.py — 2026-04-24T19:44:03.140400
% Do not edit manually. Run: uv run python scripts/update_counts.py
\newcommand{\totaltheorems}{3993}
\newcommand{\substantivetheorems}{3882}
\newcommand{\placeholdertheorems}{111}
\newcommand{\axiomcount}{1}
\newcommand{\sorrycount}{0}
\newcommand{\sorrypercent}{0.0}
\newcommand{\leanmodules}{167}
\newcommand{\leandefinitions}{3297}
\newcommand{\pythonmodules}{68}
\newcommand{\testfiles}{59}
\newcommand{\totaltests}{2836}
\newcommand{\figurecount}{111}
\newcommand{\notebookcount}{54}
\newcommand{\papercount}{19}
\newcommand{\aristotleproved}{322}
\newcommand{\aristotleruns}{44}
% --- Paper 18 / Phase 5t per-section counts (FermiHubbardDimer.lean) ---
\newcommand{\fhdTotal}{143}
\newcommand{\fhdWaveTwo}{13}
\newcommand{\fhdWaveThree}{6}
\newcommand{\fhdWaveFour}{22}
\newcommand{\fhdWaveFive}{22}
\newcommand{\fhdWaveSixABC}{38}
\newcommand{\fhdWaveSixExtra}{15}
\newcommand{\fhdWaveSeven}{19}
\newcommand{\fhdWaveEight}{8}

%% ═══════════════════════════════════════════════════════════════════
%% \figuredeferred{id}{reason} — the declared-deferral form
%% ═══════════════════════════════════════════════════════════════════
%%
%% WAVE_EXECUTION_PIPELINE.md §Stage 9: "A figure that is planned but not yet
%% drawn is declared with \figuredeferred{id}{reason}, never omitted silently."
%% The law mandated this form and, measured 2026-08-09, it had ZERO uses across
%% all 21 bundle drafts — nine of which carried no figures at all. Nine bundles
%% with no plan is a charter defect; nine bundles with a *stated* plan and a
%% stated blocker is an honest manuscript in progress, and a referee can tell
%% the two apart only if the plan is on the page.
%%
%% The deferral is DELIBERATELY VISIBLE in the compiled PDF. A silent macro
%% would satisfy `bundle_figure_adequacy` while showing a reader nothing, which
%% is the vacuous-Stage-9 failure ("ALL figures PASS" over an empty set) rebuilt
%% one layer up.
%%
%% ⚠️ The reason must name what BLOCKS the figure, not that it is absent.
%% "not yet drawn" is boilerplate and defeats the purpose; "awaiting the L=8
%% m→0 production run" is a blocker a reader can check and a later pass can
%% discharge.
%%
%% Definition lives here, beside docs/counts.tex, because 21 drafts consume it
%% and a per-draft copy is 21 places for the rendering to drift.

%% Rendered as a non-floating box on purpose. A `figure` float would be subject
%% to revtex's two-column placement rules across eleven drafts for no gain, and
%% a deferral does not need a figure number: it is a promise, not a result.

\providecommand{\figuredeferred}[2]{%
  \par\medskip\noindent
  \fbox{\parbox{0.92\columnwidth}{%
    \small\textbf{Figure deferred: \texttt{#1}.}
    Planned for this section; not yet produced. \emph{Blocked on:} #2.%
  }}%
  \par\medskip
}


\begin{document}

\title{Minimal Hypotheses for the Aliferis--Gottesman--Preskill
Threshold Recursion, Machine-Checked in Lean~4}

\author{John G.\ Roehm}
\affiliation{NetRxn Foundation}

\date{\today}

\begin{abstract}
The accuracy-threshold recursion of Aliferis, Gottesman and Preskill for the
concatenated Steane $[[7,1,3]]$ code holds under strictly weaker hypotheses than
its published proof assumes: the joint failure of a malignant pair of circuit
locations need only be sub-multiplicative, not independent, and each extended
rectangle's malignant-pair set need only be bounded in cardinality, not fixed at
it. Both weakenings matter for the models the argument is supposed to cover.
Independence is an equality between measures and therefore excludes every
negatively correlated noise model, while local stochastic noise supplies
sub-multiplicativity in general; and equality on the count is satisfied by no
real extended rectangle, whose malignant set is a proper subset of the pairs
available to it. We obtain both by deriving the recursion, rather than positing
it, from a single inclusion of failure events on a probability space: a
level-$(L{+}1)$ rectangle fails only when some malignant pair of its level-$L$
constituents both fail. Induction on that inclusion turns the
double-exponential bound into a statement about a failure \emph{probability}
rather than about a real sequence defined to satisfy its own recursion. At the
malignant-pair count Aliferis, Gottesman and Preskill report for the Steane CNOT
extended rectangle the chain discharges a numerical certificate exceeding
$2.73 \times 10^{-5}$. Two things are not derived and are stated as such
throughout: the malignant-pair count itself, which is a combinatorial
enumeration over a concrete circuit and enters as a literal, and a noise model
with a strictly positive base rate, without which the satisfiability witness
shows the hypotheses are consistent but not that the bound is ever approached.
The development is closed under the Lean kernel's standard axioms alone.
\end{abstract}

\maketitle

%% ════════════════════════════════════════════════════════════════════
\section{Introduction}
\label{sec:intro}

A fault-tolerance threshold theorem is the statement that every resource count
in quantum computing presupposes. Gate counts, qubit budgets and code-distance
estimates are all costs paid inside a regime in which arbitrarily reliable
logical computation is reachable at all, and it is the threshold theorem that
establishes there is such a regime. Its content is a recursion: below a critical
physical error rate, concatenating a code with itself suppresses the logical
error rate faster than the overhead grows, so the level-$L$ logical error rate
falls double-exponentially in $L$.

Formal verification of quantum computation is an active field, and it has not
targeted that recursion. VOQC/SQIR in Coq~\cite{Hietala2020VOQC} verifies
equivalence-preserving optimizations of quantum circuits: a certified compiler
pass guaranteeing that an optimized circuit denotes the same unitary as its
input. Qbricks~\cite{Chareton2021Qbricks} provides a deductive, Why3-based
framework for parametrized circuit-building quantum programs, with certified
case studies including quantum phase estimation and Grover search. The broader
landscape of quantum Hoare logics and tool support across Coq, Isabelle, Why3
and dedicated systems is surveyed by Lewis, Soudjani and
Zuliani~\cite{Lewis2021VerifQC}. Each of these certifies that a given circuit or
program computes what it should. None of them formalizes the argument that
makes running such a circuit reliably possible.

This paper reports what happened when we did. The target is the concatenated
Steane $[[7,1,3]]$ code under abstract local stochastic noise, following
Aliferis, Gottesman and Preskill~\cite{AGP2006} (henceforth AGP), whose threshold
argument is the one with a rigorously proven numerical bound attached. The
formalization is in Lean~4 against Mathlib \texttt{\mathlibrev}
(toolchain~\leantoolchain), and the recursion is derived from the extended
rectangle's structure rather than instantiated.

\paragraph{What is new.} Reproducing a 2006 theorem in a proof assistant is not
by itself a result, and we do not present it as one. What the formalization
produced that the prose proof does not contain is a weakening of the recursion's
hypotheses, in two places, each of which changes which physical models the bound
covers.

The first is independence. AGP's single-level step estimates the probability
that both locations of a malignant pair fail, and the natural formalization of
``distinct locations fail independently'' is Mathlib's \texttt{IndepSet}, which
asserts $\mu(s \cap t) = \mu(s)\,\mu(t)$. That is an equality, and the union
bound consumes only the inequality $\mu(s \cap t) \le \mu(s)\,\mu(t)$. Stating
the hypothesis as independence therefore excludes every negatively correlated
noise model --- models in which one location's failure makes another's less
likely --- and those lie inside the local stochastic class the theorem is
supposed to cover, not outside it. Section~\ref{sec:single} isolates
independence into a single lemma and states the recursion step from
sub-multiplicativity alone.

The second is the malignant-pair count. Written as $|F| = A$, the hypothesis
that a rectangle's malignant-pair set has cardinality $A$ is satisfied by no
extended rectangle that occurs, because a malignant set is a proper subset of
the $\binom{M}{2}$ pairs its $M$ locations admit. The bound is monotone in the
cardinality, so $|F| \le A$ is both what the proof uses and what an instance can
supply. Section~\ref{sec:levels} carries the composition in that form.

Neither weakening is an artifact of Lean. Both are statements about which noise
models the AGP recursion applies to, and both are the kind of thing a mechanized
proof surfaces because the proof assistant will not let an unused strength go
unnoticed.

\paragraph{Priority, and its limits.} To our knowledge the recursive threshold
argument itself --- the concatenation recursion carried to the
double-exponential bound with a numerical threshold attached --- has not
previously been machine-checked. Two limits bound that claim and we state them
here rather than let a referee find them. The survey behind it covered the Coq
stabilizer and verified-compilation line, the Isabelle/HOL quantum Hoare-logic
and concentration-inequality developments, and the Lean quantum-information
libraries; the Isabelle Archive of Formal Proofs was not exhaustively searched,
so the negative result there should be read as probable rather than established.
And the claim is about the threshold \emph{argument}, not its ingredients:
concentration inequalities, stabilizer formalisms and code definitions all exist
in those systems, which is precisely why the recursion that connects them is the
part worth reporting.

%% ════════════════════════════════════════════════════════════════════
\section{The model}
\label{sec:model}

\subsection{Noise}

A noisy quantum circuit is a set of \emph{locations}: gates, preparations,
measurements, and the idle intervals between them. Under an abstract local
stochastic noise model each location either operates ideally or fails, failure
at a single location occurs with probability at most $\varepsilon$, and the
failure events at distinct locations satisfy a joint-probability bound. The
model is abstract in that it says nothing about what a failure does to the
state; it is the combinatorics of which locations fail that the threshold
argument reasons about.

We are explicit at the outset that this is the only noise model treated here.
A topological error model, in which errors live on a lattice and are suppressed
by a gap rather than by a per-location rate, is a different object. It is not
this one up to constants and the two are not additive, so its published
threshold figures live on a different operational domain. Specializing to it is
deferred rather than claimed.

\subsection{Rectangles, and what makes one bad}

Concatenation implements a level-$L$ gate out of level-$(L{-}1)$ gates. The
object the argument reasons about is the \emph{extended rectangle}: a level-$L$
gate together with the error-correction rounds that flank it,
\begin{equation}
  \text{1-EC} \circ \text{1-Ga} \circ \text{1-EC}.
  \label{eq:exrec}
\end{equation}
An extended rectangle is \emph{good} when the faults inside it are few enough
that the level-$L$ gate it simulates is correct, and \emph{bad} otherwise. For a
distance-three code a single fault is always tolerable, so the leading failure
mode is two faults, and the combinatorial content of the argument is which
\emph{pairs} of locations are jointly intolerable. A pair of locations is
\emph{malignant} when some pattern of faults on that pair leaves the rectangle
incorrect; AGP's terminology is that a set of locations is benign when the
rectangle is correct for any faults on those locations, and malignant when it is
not.

Write $F$ for the malignant-pair set of a rectangle and $A$ for a bound on its
cardinality. If every malignant pair fails jointly with probability at most
$\varepsilon^2$, the rectangle is bad with probability at most $A\varepsilon^2$,
which is the level recursion
\begin{equation}
  \varepsilon_{L+1} \;\le\; A\,\varepsilon_L^{\,2}.
  \label{eq:agp-recursion}
\end{equation}
Below $\varepsilon_0 = 1/A$ the right-hand side contracts, and iterating gives
double-exponential suppression.

\subsection{What a threshold theorem asserts, and what it does not}

Equation~\eqref{eq:agp-recursion} asserts a bound on a per-level failure
probability under a per-location failure rate. It does not assert anything about
a physical device. The identification of a circuit location with a device
operation, and of a per-location failure probability with a measured error rate,
is a modelling step outside the theorem, and the threshold is a lower bound on
what a device must beat rather than a prediction about any particular hardware.
Sections~\ref{sec:single} and~\ref{sec:levels} prove the recursion and its
consequence; Section~\ref{sec:steane} supplies the count that turns it into a
number; Section~\ref{sec:scope} states what remains outside all of it.

%% ════════════════════════════════════════════════════════════════════
\section{The single-level bound, derived}
\label{sec:single}

The recursion Eq.~\eqref{eq:agp-recursion} is an inequality between
probabilities, and this section derives it as one. Fix a measurable space
$\Omega$, a measure $\mu$ on it, an index type $\iota$ of locations, and a family
$\mathrm{fail} : \iota \to \mathcal{P}(\Omega)$ of failure events.

\subsection{The union bound}

The rectangle is bad exactly when some malignant pair fails jointly, so its
failure event is the union over the malignant-pair set. The estimate on that
union uses no independence at all.

\begin{lemma}[\texttt{measure\_malignant\_union\_le}]
\label{lem:union}
Let $F$ be a finite set of pairs of locations and $\varepsilon$ an extended
non-negative real. If $\mu(\mathrm{fail}_i \cap \mathrm{fail}_j) \le
\varepsilon^2$ for every $(i,j) \in F$, then
\begin{equation}
  \mu\!\left(\bigcup_{(i,j)\in F} \mathrm{fail}_i \cap \mathrm{fail}_j\right)
    \;\le\; |F| \cdot \varepsilon^2.
  \label{eq:union-bound}
\end{equation}
\end{lemma}

The proof is the finite union bound followed by a constant sum: the measure of
the union is at most $\sum_{p \in F} \mu(\mathrm{fail}_{p_1} \cap
\mathrm{fail}_{p_2})$, each summand is at most $\varepsilon^2$ by hypothesis, and
a constant sum over $F$ is $|F|\,\varepsilon^2$. It is stated over
$\overline{\mathbb{R}}_{\ge 0}$, where no finiteness side condition is needed.

\subsection{Where independence enters, and how little of it is needed}

Lemma~\ref{lem:union} takes the per-pair estimate as a hypothesis. Two lemmas
supply it, and the difference between them is this paper's first substantive
point.

\begin{lemma}[\texttt{malignant\_pair\_indep\_le}]
\label{lem:indep}
If $\mathrm{fail}_i$ and $\mathrm{fail}_j$ are independent under $\mu$ and each
has measure at most $\varepsilon$, then $\mu(\mathrm{fail}_i \cap
\mathrm{fail}_j) \le \varepsilon^2$.
\end{lemma}

This is the only place in the development where independence is used, and its
proof is one rewrite: independence turns the measure of the intersection into a
product, and the product of two quantities each at most $\varepsilon$ is at most
$\varepsilon^2$. Isolating it makes visible that the rewrite is an equality and
the conclusion an inequality, so the hypothesis is stronger than the proof
consumes. The general form drops it:

\begin{lemma}[\texttt{exRecFailure\_le\_agpRecursionStep\_of\_submul}]
\label{lem:submul}
Suppose $\mu(\mathrm{fail}_i) \le \varepsilon$ for every location $i$, and that
every pair in $F$ is sub-multiplicative, $\mu(\mathrm{fail}_i \cap
\mathrm{fail}_j) \le \mu(\mathrm{fail}_i)\,\mu(\mathrm{fail}_j)$. Then
Eq.~\eqref{eq:union-bound} holds.
\end{lemma}

Sub-multiplicativity is implied by independence, so nothing that had the stronger
hypothesis loses anything; the converse fails, and the models in the gap are the
negatively correlated ones. A noise process in which the failure of one circuit
location suppresses the failure probability of another --- shared calibration
that a fault perturbs, a control line whose failure idles a neighbouring gate,
any mechanism that anticorrelates two locations --- satisfies
Lemma~\ref{lem:submul} and violates Lemma~\ref{lem:indep}. AGP's local
stochastic class is defined by an upper bound on joint failure probabilities,
which is sub-multiplicativity, not by an equality. The formalization does not
improve on AGP's intent; it makes the statement match it.

\subsection{From a bound to a bounded probability}

The real-valued expression $A\varepsilon^2$ is a definition, and a definition
called a union bound is not one until some probability is shown to satisfy it.
The theorem that connects them is
\texttt{exRecFailureProb\_le\_exRecFailureBound}: for a probability measure
$\mu$, a rectangle whose malignant-pair attestation carries the count $A$, a
malignant-pair set $F$ with $|F| = A$, a per-location bound $\varepsilon$ and
independence within each malignant pair, the real number
$\mu\!\left(\bigcup_{(i,j) \in F} \mathrm{fail}_i \cap \mathrm{fail}_j\right)$ is
at most $A\varepsilon^2$. Its content is the coercion: the combinatorial layer
carries location counts and malignant-pair counts and mentions no measure, and
this statement is the bridge across which the counting layer acquires
probabilistic meaning.

Two monotonicity lemmas travel with it and are used rather than quoted.
\texttt{agp\_bound\_mono\_card} states that $n \le m$ implies
$n\varepsilon^2 \le m\varepsilon^2$, which is the form a caller holding
$|F| \le A$ needs; \texttt{agp\_bound\_mono\_malignant\_set} states monotonicity
in the set, $F \subseteq G$, and is derived from the cardinality form so the two
cannot disagree.

%% ════════════════════════════════════════════════════════════════════
\section{Composition across levels}
\label{sec:levels}

\subsection{The concatenation picture as one inclusion}

Section~\ref{sec:single} bounds a single rectangle. Concatenation is what makes
the bound recursive, and the whole of it enters the formalization as one
hypothesis:
\begin{equation}
  \mathrm{fail}^{(L+1)}_j \;\subseteq\;
  \bigcup_{p \in \mathrm{malignant}(L,j)}
    \mathrm{fail}^{(L)}_{p_1} \cap \mathrm{fail}^{(L)}_{p_2}.
  \label{eq:hstruct}
\end{equation}
A level-$(L{+}1)$ extended rectangle fails only when some malignant pair of its
level-$L$ constituents both fail. That is the AGP concatenation picture written
in the language of the measure space, and compressing it to a single inclusion
of events is what the formalization forced: everything else in the composition is
Lemma~\ref{lem:submul} applied at each level with the previous level's rate.

\subsection{The induction}

Let $\mathrm{agpSeq}(A,\varepsilon_0)$ denote the sequence
$\mathrm{agpSeq}_0 = \varepsilon_0$, $\mathrm{agpSeq}_{L+1} = A\,
\mathrm{agpSeq}_L^{\,2}$, valued in $\mathbb{R}_{\ge 0}$. The choice of
$\mathbb{R}_{\ge 0}$ rather than $\mathbb{R}$ or
$\overline{\mathbb{R}}_{\ge 0}$ is what lets one induction serve both the
measure-theoretic and the real-valued sides: it embeds into
$\overline{\mathbb{R}}_{\ge 0}$ as a semiring and into $\mathbb{R}$ as the
sequence the threshold chain already reasons about, carrying neither a
non-negativity nor a finiteness side condition at each step.

\begin{theorem}[\texttt{levelFailure\_le\_agpSeq}]
\label{thm:levels}
Let $\mathrm{fail}$ be a family of failure events indexed by level and location,
$\mathrm{malignant}$ a family of finite malignant-pair sets, $A$ a natural number
and $\varepsilon_0$ a base rate. Assume
\begin{enumerate}
\item[(H1)] $\mu(\mathrm{fail}^{(0)}_i) \le \varepsilon_0$ for every level-zero
  location;
\item[(H2)] $|\mathrm{malignant}(L,j)| \le A$ for every level and rectangle;
\item[(H3)] the inclusion Eq.~\eqref{eq:hstruct};
\item[(H4)] sub-multiplicativity within each malignant pair.
\end{enumerate}
Then $\mu(\mathrm{fail}^{(L)}_i) \le \mathrm{agpSeq}(A,\varepsilon_0)_L$ for
every level $L$ and every level-$L$ rectangle $i$.
\end{theorem}

The induction is short because each ingredient is already in place: monotonicity
of the measure across Eq.~\eqref{eq:hstruct}, then Lemma~\ref{lem:submul} with
the inductive hypothesis as the per-location rate, then
\texttt{agp\_bound\_mono\_card} to pass from $|\mathrm{malignant}(L,j)| \le A$
to $A$. Hypotheses (H2) and (H4) are the weakest forms the proof uses, which is
the point of Section~\ref{sec:intro}'s second paragraph on what is new.

\subsection{The double exponential, and what it is a fact about}

The closed form of Eq.~\eqref{eq:agp-recursion} is a statement about
non-negative real sequences and has no quantum content whatever. If
$\varepsilon : \mathbb{N} \to \mathbb{R}$ is non-negative and satisfies
$\varepsilon_{L+1} \le A\,\varepsilon_L^{\,2}$ with $A \ge 0$, then
\begin{equation}
  A\,\varepsilon_L \;\le\; \left(A\,\varepsilon_0\right)^{2^L}
  \label{eq:double-exp}
\end{equation}
for every $L$ (\texttt{agp\_double\_exp\_bound}), by an induction whose step is
$A\varepsilon_{L+1} \le A(A\varepsilon_L^2) = (A\varepsilon_L)^2 \le
((A\varepsilon_0)^{2^L})^2 = (A\varepsilon_0)^{2^{L+1}}$. When
$A\varepsilon_0 < 1$ the right-hand side is below one at every level
$L \ge 1$ (\texttt{agp\_double\_exp\_bound\_lt\_one}). We state
Eq.~\eqref{eq:double-exp} flatly as standard and claim no credit for it: it is
where the double exponential comes from, and it is not a result of this paper.

What is a result of this paper is that the sequence
Eq.~\eqref{eq:double-exp} is applied to is now known to bound a probability.
\texttt{agpSeq\_coe\_eq\_agpLevelSequence} identifies the
$\mathbb{R}_{\ge 0}$-valued sequence of Theorem~\ref{thm:levels} with the
real-valued sequence the threshold chain reasons about, and the two capstones
follow. \texttt{concatenated\_failure\_double\_exp} gives
\begin{equation}
  A \cdot \mathbb{P}\!\left[\text{level-}L\text{ rectangle fails}\right]
    \;\le\; \left(A\,\varepsilon_0\right)^{2^L},
  \label{eq:capstone-doubleexp}
\end{equation}
and \texttt{steane\_concatenated\_failure\_below\_threshold} gives, below the
threshold of Section~\ref{sec:steane}, that
$A_{\mathrm{CNOT}} \cdot \mathbb{P}[\text{level-}L\text{ rectangle fails}] < 1$
for every $L \ge 1$. The distinction between Eq.~\eqref{eq:capstone-doubleexp}
and Eq.~\eqref{eq:double-exp} is the entire point of the composition: the
latter is about a sequence defined to satisfy its own recursion, and the former
is about a measure of a set of noise realizations. Their dependency closures
differ accordingly --- Eq.~\eqref{eq:capstone-doubleexp} reaches both the
measure-theoretic union bound of Section~\ref{sec:single} and the real-valued
threshold chain of Section~\ref{sec:steane}, and that joint reach is what
``derived from the extended-rectangle structure'' means operationally.

\figuredeferred{fig:d6-hypothesis-dependency}{the figure must be generated
through \texttt{src/core/visualizations.py} so the structural figure checks of
Stage~9 can see it, and no generator for a proof-dependency diagram exists in
that module yet; the content is fixed --- which of (H1)--(H4) each capstone
consumes, and which two are modelling assumptions rather than consequences}

\subsection{Satisfiability, and exactly what it establishes}

Theorem~\ref{thm:levels} and its corollaries are conditional on four
hypotheses, and a conditional theorem whose hypotheses cannot hold together
proves nothing. \texttt{levelFailure\_hypotheses\_satisfiable} exhibits, for an
arbitrary $A$ and an arbitrary base rate, a level structure satisfying all four
--- and additionally off-diagonality of the malignant pairs, which is strictly
more than any of them requires. Off-diagonality is not cosmetic: a diagonal pair
$(i,i)$ is not a pair of two distinct locations at all, AGP's malignant set
counts pairs of different locations, and under an independence hypothesis
$\mathrm{IndepSet}(s,s)$ forces $\mu(s) = \mu(s)^2$, so a diagonal malignant set
could be satisfied only where the theorem says nothing.

The witness takes every failure event empty. That is enough to rule out
vacuity and it is not enough for anything else, and we state the consequence
rather than imply it: with $\mathrm{fail} \equiv \emptyset$ every probability is
zero, so (H1) is satisfied without being exercised. The witness shows the
hypothesis set is consistent; it does not exhibit a model in which the bound is
approached. A model with a strictly positive base rate needs a genuine
independent product space over the circuit's locations, which is the same
construction as deriving the malignant-pair count itself.
Section~\ref{sec:scope} names it as open work.

%% ════════════════════════════════════════════════════════════════════
\section{The Steane instance and the threshold number}
\label{sec:steane}

\subsection{The extended rectangle and its count}

AGP instantiate the argument on Steane's $[[7,1,3]]$ code, the smallest
distance-three CSS code. Their level-1 error-correction gadget contains 142
locations --- four encoders at 18 locations each, four 1-measurements at 7, two
rests at 7, and four encoded CNOT gates at 7 --- so the CNOT extended rectangle,
comprising four such gadgets and one CNOT, contains
\begin{equation}
  M_{\rm CNOT} = 575
  \label{eq:mcnot}
\end{equation}
locations, and is the largest 1-exRec in the level-1 simulation~\cite[\S8.3]{AGP2006}.
The count is contingent on the noise model in a way worth stating: it presumes
storage faults are counted, and AGP report 487 locations if they are neglected.
The symbol $M$ is ours; AGP write $\ell$ for the maximal number of locations in a
1-Rec.

Counting malignant pairs within that rectangle is a purely combinatorial
exhaustive enumeration, which AGP carried out by computer program. The result
they report is
\begin{equation}
  A_{\rm CNOT} = 35\,235
  \label{eq:acnot}
\end{equation}
malignant pairs, obtained as the sum of the entries of a lower-triangular
$7 \times 7$ matrix in~\cite[\S8.3]{AGP2006}, immediately following that
section's Eq.~(32). Like Eq.~\eqref{eq:mcnot} it is model-contingent, and the
contingency is large: the same section gives $A_{\rm CNOT} = 22\,701$ if storage
faults are neglected and $A_{\rm CNOT} \approx 7\,183$ for depolarizing noise
rather than adversarial independent stochastic noise. Eq.~\eqref{eq:acnot} is
the adversarial figure, which is the conservative one and the one matching the
model of Section~\ref{sec:model}. The choice is not free and it moves the
threshold by a factor of five.

Its well-formedness is checked rather than asserted:
\texttt{steaneMalignancyCounts\_A\_CNOT\_le\_choose\_two} proves
$A_{\rm CNOT} \le \binom{M_{\rm CNOT}}{2} = 165\,025$, so the count cannot exceed
the number of pairs available to it.

\subsection{What the certificate certifies}

The formalized threshold is the inverse of the malignant-pair count,
\begin{equation}
  \varepsilon_0 \;:=\; \frac{1}{A_{\rm CNOT}},
  \label{eq:threshold-def}
\end{equation}
and \texttt{agp\_threshold\_steane} is Eq.~\eqref{eq:double-exp} at
$A = A_{\rm CNOT}$, for every level and every non-negative rate.
\texttt{agp\_threshold\_steane\_strict} is the operationally useful form:
strictly below Eq.~\eqref{eq:threshold-def}, every level $L \ge 1$ satisfies
$A_{\rm CNOT}\varepsilon_L < 1$, so concatenation suppresses without bound. The
numerical certificate \texttt{steaneAGPThreshold\_gt} is discharged by
\texttt{norm\_num} rather than quoted,
\begin{equation}
  \frac{1}{A_{\rm CNOT}} \;=\; 2.838 \times 10^{-5} \;>\; 2.73 \times 10^{-5},
  \label{eq:threshold-numeric}
\end{equation}
and \texttt{agp\_threshold\_steane\_numerical} bundles it with
$\varepsilon_0 \le 1$ so the quantity is pinned on both sides.

\paragraph{What Eq.~\eqref{eq:threshold-def} is, and what AGP's threshold is.}
These are not the same quantity, and conflating them would overstate the result.
AGP's published rigorous threshold is not $1/A_{\rm CNOT}$; it is the inverse of
a twice-corrected count, and the corrections go the wrong way for us. Their
chain is
\begin{equation}
  A_{\rm CNOT} = 35\,235
  \;\longrightarrow\; A'_{\rm CNOT} \approx 36\,108
  \;\longrightarrow\; A''_{\rm CNOT} \approx 36\,511,
  \label{eq:agp-chain}
\end{equation}
with $\varepsilon_0^{\rm AGP} = 1/A''_{\rm CNOT} \ge 2.739 \times 10^{-5}$,
which is~\cite[Eq.~(36)]{AGP2006} and rounds down to the
$2.73 \times 10^{-5}$ of their Theorem~3. The first correction absorbs the
\emph{cubic} term of their single-level bound: AGP's estimate is
$A_{\rm CNOT}\varepsilon^2 + B_{\rm CNOT}\varepsilon^3$ with
$B_{\rm CNOT} = \binom{575}{3}$, and $A'$ is the constant for which the pure
quadratic form dominates it. The second is a factor
$(1 - C/A'_{\rm CNOT})^{-8} \approx 1.0111$, where $C = 50$ is the number of
locations in the ancilla encoding-and-verification circuit, conditioning on
acceptance of the eight independently prepared ancilla blocks used in the four
error-correction gadgets.

The consequence for this paper is a scope statement, not a correction to the
Lean. What Section~\ref{sec:single} derives is a bound on the probability that
some malignant \emph{pair} fails, and $A_{\rm CNOT}$ is exactly the count of
those pairs, so Eq.~\eqref{eq:threshold-def} is the right constant for the
quantity actually bounded. The triple-fault contribution AGP's $B$ term carries,
and the ancilla-acceptance conditioning, are not modelled here at all. A reader
should therefore take Eq.~\eqref{eq:threshold-numeric} as certifying the
threshold of the pair-failure recursion, and should take $2.739 \times 10^{-5}$,
not $2.838 \times 10^{-5}$, as AGP's rigorous threshold for the full argument.
Extending the derivation to the cubic term is arithmetic on the same chain
rather than a new argument, and it is not done here.

\paragraph{On the number that circulates alongside it.} A threshold near
$10^{-4}$ is frequently quoted for concatenated distance-three codes. It is a
heuristic estimate, not a rigorous bound, and AGP are explicit that their own
depolarizing-noise figure of $9.376 \times 10^{-5}$ is likewise ``not a rigorous
lower bound on the accuracy threshold''. Eq.~\eqref{eq:threshold-numeric} and the
AGP chain Eq.~\eqref{eq:agp-chain} are rigorous and roughly a factor of four
smaller, and the gap between the two is a statement about proof technique rather
than about hardware.

\figuredeferred{fig:d6-threshold-double-exp}{double-exponential suppression
against level, with the malignant-pair count swept over the four values AGP
report --- $7\,183$, $22\,701$, $35\,235$, $36\,511$ --- so the model-choice
sensitivity of the preceding paragraphs is visible rather than asserted; blocked
on a generator in \texttt{src/core/visualizations.py}, which is the canonical and
only home for figure functions, and the curve is closed-form so the sweep is
cheap once the generator exists}

%% ════════════════════════════════════════════════════════════════════
\section{Scope and outlook}
\label{sec:scope}

Four limits bound this result, and none of them is implicit in the Lean.

\emph{The count is a literal.} Equation~\eqref{eq:acnot} enters from
Ref.~\cite{AGP2006} and is not derived. Deriving it means enumerating malignant
pairs over a concrete 575-location circuit, which is a combinatorial undertaking
of a different kind from anything above.

\emph{The noise model is abstract local stochastic only.} The topological error
model is a strictly different object, not equivalent to this one up to constants
and not additive, and its specialization is deferred rather than claimed.

\emph{The count is conservative rather than optimal.} A smaller malignant-pair
count raises the certified threshold, and tightening it is arithmetic on the
same chain rather than a new argument. So is extending the single-level bound to
AGP's cubic term, as Section~\ref{sec:steane} records.

\emph{The threshold is not a statement about hardware.} It is a lower bound on
what a device must beat under this model, and the identification of a circuit
location with a device operation is a modelling step that is not part of the
theorem.

The development is closed under the Lean kernel's standard axioms alone,
$\{\texttt{propext}, \texttt{Classical.choice}, \texttt{Quot.sound}\}$, with no
project-local axioms, no \texttt{sorry} and no \texttt{native\_decide}.

The open problem is the enumeration, and it is worth naming precisely because it
closes two gaps at once. Enumerating the malignant pairs of the Steane CNOT
extended rectangle means deciding, for each of the $\binom{575}{2} = 165\,025$
candidate pairs, whether some fault pattern on that pair leaves the rectangle
incorrect --- a decidable question about a fixed circuit, and the computation AGP
performed outside a proof assistant. Carrying it inside one would replace
Eq.~\eqref{eq:acnot} with a derivation, and the concrete product space it
requires over the circuit's locations is precisely the model with a strictly
positive base rate that Section~\ref{sec:levels}'s satisfiability witness does
not supply. Until then the recursion is machine-checked and the constant it is
evaluated at is quoted, and this paper is a step rather than a terminus.

%% ════════════════════════════════════════════════════════════════════
\appendix
\section{The Lean development}
\label{app:lean}

The threshold chain is nine modules under
\verb|lean/SKEFTHawking/FaultTolerance/|. Its capstone
\verb|steane_concatenated_failure_below_threshold| has a project-internal
dependency closure of 46 declarations, of which 14 are theorems, computed from
the extracted dependency graph \verb|lean/lean_deps.json| rather than by reading
imports.

\begin{table}[htbp]
\caption{Modules in the capstone's dependency closure, and what each contributes.}
\label{tab:modules}
\begin{ruledtabular}
\begin{tabular}{ll}
Module & Contributes \\
\colrule
\verb|Basic| & shared definitions \\
\verb|StabilizerCode| & stabilizer substrate \\
\verb|SteaneCode| & the $[[7,1,3]]$ instance \\
\verb|Counting| & Eqs.~\eqref{eq:mcnot},~\eqref{eq:acnot}; well-formedness \\
\verb|MalignantUnionBound| & Lemmas~\ref{lem:union}--\ref{lem:submul} \\
\verb|Concatenation| & the level sequence \\
\verb|DoubleExp| & Eq.~\eqref{eq:double-exp} \\
\verb|ConcatenatedComposition| & Theorem~\ref{thm:levels}; the capstones \\
\verb|AGP/Threshold| & Eqs.~\eqref{eq:threshold-def},~\eqref{eq:threshold-numeric} \\
\end{tabular}
\end{ruledtabular}
\end{table}

Four further modules in the same directory --- \verb|ExRec|, \verb|Malignant|,
\verb|NoiseModel| and \verb|NoiseModelMT| --- formalize the objects the argument
reasons about and are reached by the capstone only through
\verb|MalignantUnionBound|'s bridge to \verb|exRecFailureBound|. We report that
rather than presenting them as part of the proved chain.

Build provenance: Lean toolchain \leantoolchain, Mathlib revision
\mathlibrev.

%% ════════════════════════════════════════════════════════════════════
\begin{thebibliography}{99}

\bibitem{AGP2006}
P.~Aliferis, D.~Gottesman, and J.~Preskill,
``Quantum accuracy threshold for concatenated distance-3 codes,''
Quantum Inf.\ Comput.\ \textbf{6}, 97--165 (2006)
[arXiv:quant-ph/0504218].

\bibitem{Hietala2020VOQC}
K.~Hietala, R.~Rand, S.-H.~Hung, X.~Wu, and M.~Hicks,
``A verified optimizer for quantum circuits,''
Proc.\ ACM Program.\ Lang.\ \textbf{5} (POPL), Article 37 (2021)
[arXiv:1912.02250], doi:10.1145/3434318.

\bibitem{Chareton2021Qbricks}
C.~Chareton, S.~Bardin, F.~Bobot, V.~Perrelle, and B.~Valiron,
``An automated deductive verification framework for circuit-building
quantum programs,''
in \emph{Programming Languages and Systems (ESOP 2021)},
Lecture Notes in Computer Science \textbf{12648}, 148--177
(Springer, 2021) [arXiv:2003.05841], doi:10.1007/978-3-030-72019-3\_6.

\bibitem{Lewis2021VerifQC}
M.~Lewis, S.~Soudjani, and P.~Zuliani,
``Formal verification of quantum programs: theory, tools, and
challenges,''
ACM Trans.\ Quantum Comput.\ \textbf{5}, 1 (2024)
[arXiv:2110.01320], doi:10.1145/3624483.

\end{thebibliography}

\end{document}
